% Problem record op_64046727b81b4024. This TeX file is derived from
% database/problems_json/op_64046727b81b4024.json by scripts/sync-tex.mjs; edit the JSON record
% and rerun the script rather than editing this file.
\section{Universality of every fixed non-Gaussian polynomial generator}
\label{sec:op_64046727b81b4024}

\paragraph{Problem.}
Does every fixed non-Gaussian polynomial Hamiltonian become universal when all Gaussian controls are available?
Fix $n\geq1$.
On $L^2(\mathbb R^n)$, let $q_j,p_j$ satisfy $[q_j,p_k]=i\delta_{jk}$.
Let $H_*$ be any real Weyl-ordered polynomial in these operators of degree greater than two.
Assume that $H_*$ is essentially self-adjoint on the Schwartz space $\mathcal S(\mathbb R^n)$ of smooth, rapidly decreasing functions.
Use its unique self-adjoint closure to define $e^{-itH_*}$.
Allowed gates are $e^{-itH_*}$ for arbitrary real $t$, together with all Gaussian unitaries.
Gaussian unitaries are generated by real polynomials of degree at most two in the canonical operators.

Let $U$ be any unitary satisfying $U\mathcal S(\mathbb R^n)\subseteq\mathcal S(\mathbb R^n)$.
Fix finite $E\geq0$ and $\varepsilon>0$.
Does some finite product $V$ of allowed gates satisfy
\begin{equation}
 \sup_{\rho_{AR}:\operatorname{Tr}(\rho_A N)\leq E}
 \left\|[(\mathcal U-\mathcal V)\otimes\operatorname{id}_R](\rho_{AR})\right\|_1
 <\varepsilon,
 \qquad N=\frac12\sum_{j=1}^n(q_j^2+p_j^2-1)?
 \label{eq:fixed-poly-target}
\end{equation}
In Eq.~\eqref{eq:fixed-poly-target}, $\mathcal U(\rho)=U\rho U^\dagger$, $\mathcal V(\rho)=V\rho V^\dagger$, and $R$ is an arbitrary reference system.
The fixed $H_*$ must work for every target $U$, energy $E$, and accuracy $\varepsilon$.

\subsection*{Status}

Unsolved

\subsection*{Source}

The Discussion of Arzani, Booth, and Chabaud explicitly identifies the missing rigorous equivalence between Gaussian controls plus any fixed higher-degree polynomial and arbitrary polynomial evolutions \sourcecite{ref:fixed-poly-abc}{ABC25}.
This formulation specifies the operator domain and energy-constrained approximation target.

\subsection*{Progress}

\begin{itemize}
  \item Arzani, Booth, and Chabaud prove universality when the polynomial Hamiltonian can be chosen for each target.
Their Theorem~2 realizes arbitrary finite Fock blocks; Theorem~3 supplies unitary approximation \sourcecite{ref:fixed-poly-abc}{ABC25}.
A polynomial in total photon number can be added outside the retained block to obtain a self-adjoint realization.
These generators depend on the desired finite-dimensional action.

  \item The same paper's Theorem~4 gives a Solovay--Kitaev result for finite-block gate sets constructed from Theorem~2.
Its hypotheses do not cover every preassigned higher-degree $H_*$.
The Discussion expressly leaves the fixed-generator equivalence open \sourcecite{ref:fixed-poly-abc}{ABC25}.
\end{itemize}

\subsection*{References}

\begin{enumerate}[label={},leftmargin=4.8em,labelsep=0.5em]
  \footnotesize
  \item[\textup{[ABC25]}]\label{ref:fixed-poly-abc}
  F. Arzani, R. I. Booth, and U. Chabaud, "Effective Descriptions of Bosonic Systems Can Be Considered Complete," \emph{Nature Communications} \textbf{16}, 9744 (2025). \href{https://doi.org/10.1038/s41467-025-64872-3}{doi:10.1038/s41467-025-64872-3}; \href{https://arxiv.org/abs/2501.13857}{arXiv:2501.13857}.
\end{enumerate}

\subsection*{Comment}

The gap is a proof of the approximation in Eq.~\eqref{eq:fixed-poly-target} for every essentially self-adjoint polynomial $H_*$ of degree greater than two, or a counterexample.
Formal Lie-algebra generation alone does not justify products of unbounded evolutions or their convergence.
Arbitrary polynomial universality is therefore a related solved problem with a different resource model.

\subsection*{Field}

Quantum algorithm

\subsection*{Topic}

Quantum circuit complexity; Hamiltonian simulation

\subsection*{ID}

\texttt{op\_64046727b81b4024}
